\section{Refinement of the models}

Since the paper's first version (\code{paper-v1}), CM1 has been forced through the whole column and read against the
GCM budget by budget, defects and Earth values in both codes have been repaired, PlaSim's cold cloud has been fitted
at lunar gravity, and the scenario has gained an orientation computed from its water. The rings and the boxes forced
from the ten-level states held CM1's air toward the GCM's from 8~km, a height that leaves an Earth-scale boundary
layer and its convection free, while lunar convection reaches 55--80~km, so the hold covered most of the convecting
column and, over the Near-side Sea, took the place of the deep convection (Verification).\claim{H39} Forced instead
with the GCM's heating and moistening through the whole column, its temperature and moisture held to the GCM's only
above the GCM's convection top, CM1 convects through the depth in which the GCM convects, and the height of the hold
follows a quantity the forcing GCM computes. A coupling height fixed on Earth's scale thus covers the convection of a
world with a deeper troposphere, and one taken from the forcing model's convection top serves any world whose GCM
computes that top.\claim{R14}

Each fault found in CM1 left output that looked plausible on its own, and each was found by setting CM1 against
something outside it. Its simple boundary layer's density fault appeared when a box's boundary layer took up more
heat and water than its surface supplied; the gust over water, a formula that omits gravity, when the sea box's
surface wind was set beside the GCM's; Morrison's Earth limits on rain, in CM1's code; and the error in reading its
radiative output, which overstated CM1's radiative heating by 38--40\% in the comparisons of the first version, when
RRTMG, run offline on the same columns, failed to reproduce the cooling CM1 had written.\claim{H34}\claim{H38}\claim{N03}\claim{V08}

The forcing also shaped what CM1 showed. Under a ten-level state with Kuo-type convection, whose circulation sinks
over the Near-side Sea after midnight while the scheme keeps raining, CM1's sea box, forced through the whole column,
took the sinking and rained 0.14~mm a day after midnight; forced in the same way from a fifteen-level state with
simplified Betts--Miller convection, it rained 1.62--1.88~mm a day, against that state's own 2.28.\claim{G25} A
cloud-resolving model forced from a GCM takes on the GCM's circulation, so a difference between the two can lie in
the forcing as well as in either model's physics, and boxes that make their own large-scale ascent change the
comparison's result at every site (Results).\claim{G26} Compared at the sea surface, the two models' exchange differed
in its law as well; on fifteen levels, with its lowest level in the surface layer, PlaSim takes CM1's own similarity
theory over ice-free water with Beljaars's gust, so the two models exchange heat and vapour with the sea by one law.
What remains between them, in the timing of the rain and the temperature of the upper air, lies in the GCM's
convection and in how the boxes receive their ascent, both still under test (Next steps).\claim{N05}

PlaSim's calibrations hold near the state they were fitted on. The cold cloud's factors, computed at lunar gravity,
cooled the climate by 2.8~K and cut its rain by 19\%, the largest of the refinements measured singly, and a longwave
correction fitted on one state misplaces the cooling of another, so it is refitted whenever the state, the levels or
the convection change (Repairs).\claim{G19}\claim{R25} The dimming that holds the design temperature is then set with
whatever imbalance the physics leaves, and the design case's was set while a fault of the adaptation's own counted
the night's twilight twice: PlaSim's diagnostic of cloud forcing writes the shortwave only where the Sun is up, so the
twilight the curved-air patch adds below the horizon accumulated over its two passes, and the design case reaches its
target with too much heat on the night side and toward the poles (Verification).\claim{H40} A design climate that
holds its target temperature is therefore also checked for where its heat goes.

The scenario changed as well, with the water it adds. Angular momentum holds a body's spin axis fixed in space, so a
load that tilts the body's principal axes turns the body beneath that axis until its axis of greatest moment lies
along the spin and, under a tidal lock, its axis of least moment points to the host. As the Moon's basins fill, the
Moon follows its new axes and, while its crust responds elastically, turns by about 8\textdegree, a turn its mantle
would ease only over geological time in the central case of its present thermal state, so the runs take the elastic
state.\claim{S17} At the same sunlight the turned Moon's climate runs 0.8~K warmer, its near side's north wetter and
its southern highlands drier.\claim{G24} A world's orientation is thus an output of its scenario: an ensemble that
adds water to a body derives the orientation from the water's load before running the body's climate, and derives it
again whenever that climate's runoff moves the lakes, since the order in which the basins fill sets the path from
today's orientation to the turned one.
